% ECONOMICS_PROBLEM: {"id": "J62-562", "title": "Noncompetes, no-poach rules and worker mobility", "jel_code": "J62", "source_id": "OP-0030"}
% !TeX program = xelatex
\documentclass[11pt,a4paper]{article}
\usepackage[margin=23mm]{geometry}
\usepackage{fontspec}
\setmainfont{Latin Modern Roman}
\usepackage{amsmath,amssymb,mathtools}
\usepackage{xcolor,enumitem,booktabs,longtable,array,xurl,hyperref}
\definecolor{muted}{HTML}{53636C}
\hypersetup{colorlinks=true,urlcolor=blue,linkcolor=blue}
\setlength{\parindent}{0pt}
\setlength{\parskip}{5pt}
\newcommand{\R}{\mathbb R}
\newcommand{\E}{\mathbb E}
\newcommand{\N}{\mathbb N}
\newcommand{\ind}{\mathbf 1}
\newcommand{\Var}{\operatorname{Var}}
\newcommand{\Cov}{\operatorname{Cov}}
\newcommand{\supp}{\operatorname{supp}}
\DeclareMathOperator*{\argmin}{arg\,min}
\DeclareMathOperator*{\argmax}{arg\,max}
\newcommand{\entrymeta}[3]{{\sffamily\small\color{muted}\textbf{\texttt{#1}}\quad|\quad #2\par #3\par}}
\newcommand{\Problemref}[1]{\texttt{#1}}
\newcommand{\JELref}[2]{\texttt{#2} (JEL \texttt{#1})}
\begin{document}
\section*{Noncompetes, no-poach rules and worker mobility}
\textbf{Problem ID:} \texttt{J62-562}\quad \textbf{JEL:} \texttt{J62}\par
% BEGIN ECONOMICS STATEMENT
\label{problem:OP-0030}
\label{atlas:prob:L10}
\noindent\textbf{Original atlas ID:} L10 (entry 30). \textbf{Classification:} Editorial dynamic welfare design under mobility restrictions.

\noindent\textbf{Original atlas question:} What restrictions maximize knowledge investment without suppressing mobility and wages?

\subsection*{Definitions and assumptions}
Fix an economy, baseline workforce, and horizon $H$. Let $\ell=(\ell_c,\ell_p,\ell_d)$ specify noncompete enforcement, employer no-poach restrictions, and contract disclosure. These are different policy margins. A dynamic model $M$ includes worker matching and mobility, training and research investment, knowledge diffusion, firm entry, contract information, and product-market competition. State the feasible policy set $\mathcal L$, preferences, discount factor $\beta\in(0,1)$, and an equilibrium selection rule. Define $u_i^M(\ell,t)$ as worker consumption-and-amenity utility, $v_j^M(\ell,t)$ as owners' utility, and $c_t^M(\ell)$ as public resource costs, using announced welfare weights $\omega_i,\eta_j\geq0$ and a common utility accounting convention.
\subsection*{Formal research question}
Assuming convergence, define
\[
 W_M(\ell)=\mathbb E\sum_{t=0}^{H}\beta^t
   \left[\sum_i\omega_i u_i^M(\ell,t)
        +\sum_j\eta_j v_j^M(\ell,t)-c_t^M(\ell)\right].
\]
Characterize the set of welfare-optimal feasible restrictions, conditional on $M$, and the identified set of $W_M(\ell)-W_M(\ell_0)$ across models fitting wage, mobility, investment, startup, and innovation data. For every policy contrast, jointly estimate labor-market responses and real investment/knowledge outcomes. Determine when a mobility loss is compensated by an investment gain under the stated weights, rather than assuming that either effect necessarily dominates.
\subsection*{Interpretation and scope}
Wages and profits are distributional outcomes; they should not be added again as resource benefits if already represented by consumption utility. Patent counts are noisy innovation measures and need explicit quality or productivity validation. Enforcement changes can alter perceived constraints even where contracts are invalid, so usage, notice, beliefs, and legal enforceability must be distinguished. A historical reform identifies a particular jurisdictional contrast under its comparison assumptions; it does not reveal a universal optimal duration or geographic scope. No-poach agreements between employers have a different treatment and competitive mechanism from employee noncompetes. Evidence for one cannot establish the other's causal effects without an additional argument. Current legal validity is outside this mathematical formulation.
\subsection*{Source audit and status}
The three atlas citations concern noncompetes. Their no-poach coverage is therefore indirect and is not represented as verified no-poach evidence. [S3] has a 2021 online date and a 2022 issue date. Added [S4] provides 2026 field-experimental evidence. The dynamic welfare target is editorial; the cited mobility and wage results are already established for their settings, while no universal optimum or universal unresolved status is asserted.

\subsection*{References}
\begin{enumerate}[label=\textnormal{[S\arabic*]},leftmargin=*]
\item Matt Marx, Deborah Strumsky, and Lee Fleming (2009). \emph{Mobility, Skills, and the Michigan Non-Compete Experiment}. Management Science 55(6). \url{https://doi.org/10.1287/mnsc.1080.0985}. \textit{Locator:} Primary publisher, working-paper, or author record: bibliographic information and abstract; full-text theorem verification is not claimed. \textit{Role:} Historical policy-change evidence on inventor mobility.
\item Evan P. Starr, J. J. Prescott, and Norman D. Bishara (2021). \emph{Noncompete Agreements in the US Labor Force}. The Journal of Law and Economics 64(1), 53--84. \url{https://doi.org/10.1086/712206}. \textit{Locator:} Primary publisher, working-paper, or author record: bibliographic information and abstract; full-text theorem verification is not claimed. \textit{Role:} National survey evidence on noncompete use and employee outcomes.
\item Michael Lipsitz and Evan Starr (2022). \emph{Low-Wage Workers and the Enforceability of Noncompete Agreements}. Management Science 68(1), 143--170; published online in 2021. \url{https://doi.org/10.1287/mnsc.2020.3918}. \textit{Locator:} Primary publisher, working-paper, or author record: bibliographic information and abstract; full-text theorem verification is not claimed. \textit{Role:} Oregon policy-change evidence on low-wage workers; identifies a specific margin. The atlas 2022 issue year is consistent with volume 68(1); the publisher displays a 2021 online-publication date.
\item Bo Cowgill, Brandon Freiberg, and Evan Starr (2026). \emph{Clause and Effect: Theory and Field Experimental Evidence on Noncompete Clauses}. The Quarterly Journal of Economics, advance article qjag040. \url{https://doi.org/10.1093/qje/qjag040}. \textit{Locator:} Primary publisher, working-paper, or author record: bibliographic information and abstract; full-text theorem verification is not claimed. \textit{Role:} More recent field-experimental evidence and theory; not a universal welfare optimum.
\end{enumerate}

\subsection*{Common conventions: Source mathematical conventions}

Unless an entry states otherwise, agent and firm sets are finite. State and action spaces are measurable, strategies and outcome maps are measurable, probability laws are specified, and expectations exist for the targets under discussion. Infinite-horizon discount factors lie in $(0,1)$ and discounted objectives are finite. Norms, tolerances and complexity input sizes must be specified for computational comparisons. Constants or primitives that appear locally are inputs, not empirical facts asserted by this catalogue.

\textbf{Identification.} A statistical model $m\in\mathcal M$ implies an observable law $P_m$ and target $\theta(m)$. At law $P$, its identified set is
\[
\Theta(P;\mathcal M)=\{\theta(m):m\in\mathcal M,\ P_m=P\}.
\]
Point identification holds when this set is a singleton, or equivalently when $P_{m_1}=P_{m_2}$ implies $\theta(m_1)=\theta(m_2)$ for all compatible models. Sharp bounds mean bounds attained, or approached when the boundary is not attained, by compatible models. Writing this set is a definition, not a solution: the substantive task is to establish informative restrictions and derive or estimate the set. An unrestricted-confounding model may give only trivial bounds.

\textbf{Inference.} Identification, estimability and valid uncertainty quantification are separate questions. Fix $\alpha\in(0,1)$. A confidence procedure $C_n$ over a declared law class $\mathcal P$ has asymptotic uniform coverage at level $1-\alpha$ when
\[
\liminf_{n\to\infty}\inf_{P\in\mathcal P}\Pr_P\{\theta(P)\in C_n\}\geq1-\alpha.
\]
For set-identified targets the coverage event must be defined separately, for example coverage of the full identified set. If an entry uses identification-indexed models instead of uniquely indexed laws, the same coverage convention is applied over those models. Pointwise limits do not establish uniform validity, and asymptotic validity does not guarantee finite-sample precision. Sample sizes, dependence structures and weak-identification sequences are part of each application.

\textbf{Counterfactuals.} $Y_i(a)$ is a potential outcome under a specified intervention, including its timing, implementation and versions. It is not $Y_i$ conditional on voluntarily choosing $a$. A full intervention vector permits interference; shorthand scalar treatment notation does not silently impose no interference. Assignment, selection, attrition, anticipation and measurement must be specified. A claim about an instrument's local effect is not automatically a population policy effect.

\textbf{Economic equilibrium.} A counterfactual model includes best responses, constraints, feasibility, market clearing, state transitions and expectations consistent with the stated information. When several equilibria exist, either a declared selector is an input or the target is set-valued. Comparisons across policy regimes must use the stated selection convention. Reduced-form relationships are not presumed invariant after an equilibrium-changing intervention.

\textbf{Optimization and welfare.} Optimization is over the stated feasible set. If attainment is not established, $\sup$ or $\inf$ and $\varepsilon$-optimal choices replace an asserted maximizer or minimizer. For example, with value $V=\sup_{a\in\mathcal A}W(a)<\infty$, an $\varepsilon$-optimal set is $\{a\in\mathcal A:W(a)\geq V-\varepsilon\}$ for $\varepsilon>0$. Benefits, costs, risk and health or environmental outcomes can be aggregated only using declared units and conversion weights. Interpersonal utility weights, the treatment of fiscal transfers and foreign profits, discounting, and the behavioral welfare criterion are explicit inputs. A comparative-static derivative additionally requires existence and appropriate differentiability of the selected counterfactual.

\textbf{Sources.} Local labels [S1], [S2], and so forth restart in every question. Locators identify the relevant abstract, model, section or result at the level actually checked; a bibliographic or abstract-level verification is not represented as inspection of an entire proof. Working papers are distinguished from later publications and changing versions. Source summaries are paraphrased. All URL checks and editorial formulations refer to the audit date stated above.
% END ECONOMICS STATEMENT
\end{document}
